\subsection{Архитектура Apache Paimon}

Apache Paimon~\cite{paimon-docs} --- открытый табличный формат, объединяющий слой метаданных в стиле
Apache Iceberg с организацией данных в виде LSM-дерева~\cite{oneil1996lsm, luo2020lsm} в объектном хранилище. Ниже
рассматриваются модель данных, устройство таблицы с первичным ключом, путь чтения и
точка подключения файлового формата.

\textbf{Модель данных и метаданные.} Таблица Paimon --- это директория в объектном
хранилище, содержащая поддиректории данных и метаданных. Каждая фиксация (commit)
порождает новый \textit{снимок} (snapshot) --- небольшой файл, ссылающийся на список
\textit{манифестов}. Манифест перечисляет файлы данных, добавленные или удалённые
данной фиксацией, вместе со статистиками по столбцам каждого файла (минимум,
максимум, число null). Чтение всегда выполняется относительно конкретного снимка,
что обеспечивает изоляцию: параллельная запись создаёт новые снимки, не затрагивая
тот, который читает потребитель. Старые снимки и недостижимые файлы удаляются по
политике хранения (retention).

\textbf{Разбиение и бакеты.} Таблица может быть разбита на разделы (partitions) по
значениям выбранных столбцов; внутри раздела данные распределяются по фиксированному
числу \textit{бакетов} (buckets) хешированием ключа. Бакет --- единица параллелизма
записи и единица LSM-дерева: каждый бакет имеет собственное LSM-дерево. Запись в
бакет однопоточна, что устраняет конкуренцию за одно дерево; число бакетов задаёт
максимальный параллелизм записи.

\textbf{Таблица с первичным ключом и LSM-дерево.} В таблице с первичным ключом строки
идентифицируются ключом; обновления и удаления не перезаписывают данные на месте, а
добавляются как новые версии. Внутри бакета новые записи буферизуются в памяти в
отсортированном по ключу виде и при фиксации (по чекпойнту вычислительного движка)
сбрасываются на диск одним отсортированным неизменяемым файлом --- это \textit{уровень~0}
(\(L_0\)) LSM-дерева. Файлов \(L_0\) накапливается по одному на чекпойнт. Фоновая
операция \textit{слияния} (compaction) объединяет несколько отсортированных серий
(sorted runs) в файлы нижележащих уровней (\(L_1\), \(L_2\), \dots), отбрасывая
перезаписанные и помеченные на удаление записи и укрупняя файлы; слияние не
блокирует запись. Параметры \texttt{num-sorted-run.compaction-trigger} (число серий,
при котором запускается слияние) и \texttt{num-sorted-run.stop-trigger} (порог
обратного давления на запись) управляют его агрессивностью; \texttt{target-file-size}
задаёт целевой размер укрупнённых файлов. Схематически устройство показано на
рисунке~\ref{fig:lsm}.

\begin{figure}
\centering
\begin{tikzpicture}[
    font=\small,
    box/.style={draw, rounded corners=2pt, minimum height=7mm, minimum width=14mm, align=center},
    lvl/.style={font=\small\bfseries},
    >={Stealth[length=2mm]}
  ]
  % memory buffer
  \node[box, fill=black!5] (mem) at (0,3.4) {буфер в памяти\\(сортировка по ключу)};
  % L0
  \node[lvl] at (-4.6,2.2) {$L_0$};
  \node[box, fill=black!8] (l0a) at (-3.2,2.2) {файл};
  \node[box, fill=black!8] (l0b) at (-1.6,2.2) {файл};
  \node[box, fill=black!8] (l0c) at (0.0,2.2) {файл};
  \node at (1.4,2.2) {$\dots$};
  % L1
  \node[lvl] at (-4.6,1.0) {$L_1$};
  \node[box, fill=black!12, minimum width=22mm] (l1a) at (-2.8,1.0) {файл};
  \node[box, fill=black!12, minimum width=22mm] (l1b) at (-0.2,1.0) {файл};
  % L2
  \node[lvl] at (-4.6,-0.2) {$L_2$};
  \node[box, fill=black!16, minimum width=34mm] (l2a) at (-2.0,-0.2) {файл};
  \node[box, fill=black!16, minimum width=34mm] (l2b) at (2.0,-0.2) {файл};
  % arrows
  \draw[->] (mem) -- node[right]{\footnotesize flush по чекпойнту} (l0c);
  \draw[->] (l0a) -- ($(l1a.north)+(-0.4,0)$);
  \draw[->] (l0b) -- ($(l1a.north)+(0.4,0)$);
  \draw[->, dashed] (l0c) -- (l1b.north);
  \node[font=\footnotesize] at (3.3,1.55) {слияние ($L_0\!\to\!L_1$)};
  \draw[->] (l1a) -- (l2a.north);
  \draw[->] (l1b) -- ($(l2a.north)+(0.8,0)$);
  \node[font=\footnotesize] at (3.6,0.45) {слияние ($L_1\!\to\!L_2$)};
  % read side
  \node[box, fill=black!5, minimum width=30mm] (rd) at (5.6,3.0) {итератор слияния\\при чтении};
  \draw[->, thick] (l0c.east) to[out=0,in=180] ($(rd.west)+(0,0.3)$);
  \draw[->, thick] (l1b.east) to[out=0,in=180] (rd.west);
  \draw[->, thick] (l2b.east) to[out=0,in=180] ($(rd.west)+(0,-0.3)$);
\end{tikzpicture}
\caption{Устройство LSM-дерева бакета в Apache Paimon: буферизация в памяти, сброс в~\(L_0\) по чекпойнту, фоновое слияние на нижележащие уровни, объединение версий итератором слияния при чтении}
\label{fig:lsm}
\end{figure}

\textbf{Поле последовательности и вид изменения.} Чтобы итератор слияния выбирал
актуальную версию строки, таблица настраивается на поле последовательности
(\texttt{sequence.field}) --- монотонно возрастающую величину (например, момент
изменения), по максимуму которой определяется последняя версия ключа; поле вида
изменения (\texttt{rowkind.field}) различает вставку, обновление и удаление. Эта
схема естественно соответствует потокам захвата изменений (CDC).

\textbf{Путь чтения.} Чтение бакета порождает итератор слияния (merge iterator),
который параллельно читает отсортированные серии всех уровней и выдаёт строки в
порядке ключа, оставляя для каждого ключа версию с максимальным значением поля
последовательности. Если файлов \(L_0\) накопилось много (слияние не успело за
записью), итератор вынужден открыть и слить больше файлов --- стоимость чтения растёт.
Перед открытием файла применяется пропуск данных по статистикам из манифеста: если
диапазоны значений столбцов файла не пересекаются с предикатом запроса, файл
пропускается целиком; внутри файла дополнительно работает пропуск блоков средствами
самого файлового формата. Декодированные данные передаются вычислительному движку
пакетами строк (vectorized batches).

\textbf{Точка подключения файлового формата.} Файловый формат подключается к Paimon
через интерфейс \texttt{FileFormat} (модуль \texttt{paimon-common}), который
предоставляет фабрику записи (\texttt{FormatWriterFactory}) и фабрику чтения
(\texttt{FormatReaderFactory}). Фабрика записи создаёт по пути в хранилище объект
\texttt{FormatWriter}, принимающий строки и сериализующий их в файл выбранного
формата. Фабрика чтения создаёт по пути, проекции столбцов и (опционально)
предикату объект \texttt{FormatReaderFactory.FormatReader}, выдающий итератор пакетов
строк. Реализация формата также объявляет, какие предикаты она способна принять для
проталкивания, и формирует объект статистик по записанному файлу для манифеста.
Существующие реализации --- \texttt{paimon-format} (Parquet, ORC, Avro); добавление
модуля \texttt{paimon-vortex}, реализующего этот интерфейс поверх нативной библиотеки
Vortex, и составляет практическую часть работы.

\textbf{Интеграция с вычислительным движком.} В работе используется Apache Flink~\cite{carbone2015flink}:
потоковая запись --- через приёмник Paimon, чтение --- через источник Paimon в
пакетном режиме. Запись управляется чекпойнтами Flink: каждый чекпойнт фиксирует
снимок Paimon. Интервал чекпойнтов задаёт частоту фиксаций: чем он больше, тем
крупнее (и реже) файлы \(L_0\), но тем больше данных буферизуется и тем заметнее
всплеск ввода-вывода при сбросе. Этот параметр существенен для near-real-time-сценария
и отдельно варьируется в экспериментах.
